%============================================================
%  Bd. 09 - Auswahlprinzip
%============================================================

\documentclass[main.tex]{subfiles}

\ifSubfilesClassLoaded{
  \BandSetupStandalone{B09}
}{
}

\title{Bd. 09 - Auswahlprinzip}
\author{Martin Kunze}
\date{}
\setcounter{file}{9}

\hypersetup{
  pdftitle={Bd. 09 - Auswahlprinzip},
  pdfauthor={Martin Kunze}
}

\begin{document}

\title{Bd. 09 - Auswahlprinzip}
\author{Martin Kunze}
\date{}
\hypersetup{pageanchor=false}
\maketitle
\hypersetup{pageanchor=true}
\setcounter{tocdepth}{3}
\tableofcontents
\setcounter{file}{9}
\renewcommand{\FormulaBandID}{9}

\chapter{Einleitung}

Das Auswahlprinzip ist ein eigenständiger Axiom- und Äquivalenzkomplex. Es
wird daher nicht der elementaren Theorie surjektiver Funktionen zugeschlagen.
Dieser Band setzt die funktionsfreie Familienform aus Band~03, totale
Relationen aus Band~04 sowie die Funktionsbegriffe aus den Bänden B05--B08
voraus und entwickelt daraus Relations-, Familien- und Surjektionsform.

\chapter{Sektionen und Rechtsinverse}

\section{Rechtsinverse und Injektivität}

\FormulaThmDelta[Rechtsinverse ist injektiv]{%
  F\circ G=\Id_B\vdash G\colon B\inj A%
}{
  \DeltaRow{Mengen}{A\dsep B\dsep x\dsep y}
  \DeltaPrem{Funktionen}{F\colon A\to B \dsep G\colon B\to A}
}
\begin{tabproofwide}
  \proofstepwidestar[1]{F\circ G=\Id_B}{\rA}

  \proofstepwidestar[2]{x\in B}{\rA}
  \proofstepwidestar[3]{y\in B}{\rA}
  \proofstepwidestar[4]{G(x)=G(y)}{\rA}

  \proofstepwidestar[1]{\Id_B=F\circ G}{%
    \FormulaRefAuto{X=Y\vdash Y=X}}

  \proofstepwide[2]{x}{=}{\Id_B(x)}{%
    \FormulaRefAuto{x\in A\vdash x=\Id_A(x)}}

  \proofstepwide[1,2]{}{=}{(F\circ G)(x)}{%
    \rIE{5,2}}

  \proofstepwide[2]{}{=}{F(G(x))}{%
    \FormulaRefAuto{x\in B \vdash (F\circ G)(x)=F(G(x))}}

  \proofstepwide[2,3,4]{}{=}{F(G(y))}{%
    \rIE{4,8}}

  \proofstepwide[3]{}{=}{(F\circ G)(y)}{%
    \FormulaRefAuto{x\in B \vdash F(G(x))=(F\circ G)(x)}}

  \proofstepwide[1,3]{}{=}{\Id_B(y)}{%
    \rIE{1,3}}

  \proofstepwide[3]{}{=}{y}{%
    \FormulaRefAuto{x\in B\vdash \Id_B(x)=x}}

  \proofstepwide[1,2,3,4]{x}{=}{y}{%
    \rChain{6,12}}

  \proofstepwidestar[1]{G\colon B\inj A}{%
    \FormulaRefAuto{Injektive Funktion}{13}}
\end{tabproofwide}

\FormulaThmDelta{%
  G\colon A\sur B\dsep F\colon B\to A\dsep
  \forall y\in B\,\bigl(G(F(y))=y\bigr)\vdash
  F\colon B\inj A
}{%
  \DeltaDecl{Mengen}{A\dsep B\dsep y}%
  \DeltaDecl{Funktionen}{F\dsep G}%
}
\begin{tabproof}
  \proofstep{1}{G\colon A\sur B}{\rA}
  \proofstep{2}{F\colon B\to A}{\rA}
  \proofstep{3}{\forall y\in B\,\bigl(G(F(y))=y\bigr)}{\rA}

  \proofstep{4}{y\in B}{\rA}

  \proofstep{2,4}{(G\circ F)(y)=G(F(y))}{%
    \FormulaRefAuto{x\in A \vdash (G\circ F)(x)=G(F(x))}{4}}

  \proofstep{3,4}{G(F(y))=y}{\rRE{4,\rUE{3}}}

  \proofstep{2,3,4}{(G\circ F)(y)=y}{%
    \FormulaRefAuto{a = b,\, b = c \vdash a = c}{5,6}}

  \proofstep{4}{\Id_B(y)=y}{%
    \FormulaRefAuto{x\in A\vdash \Id_A(x)=x}{4}}

  \proofstep{4}{y=\Id_B(y)}{%
    \FormulaRefAuto{a=b\vdash b=a}{8}}

  \proofstep{2,3,4}{(G\circ F)(y)=\Id_B(y)}{%
    \FormulaRefAuto{a = b,\, b = c \vdash a = c}{7,9}}

  \proofstep{2,3}{\forall y\in B\,\bigl((G\circ F)(y)=\Id_B(y)\bigr)}{%
    \rUI{\rRI{4,10}}}

  \proofstep{1}{G\colon A\to B}{\FormulaRefAuto{Surjektive Funktion}[def]{1}}
  \proofstep{1,2}{G\circ F\colon B\to B}{%
    \FormulaRefAuto{CompositionDef}[def]{2,12}}
  \proofstep{}{\Id_B\colon B\to B}{\FormulaRefAuto{IdA}[thm]}
  \proofstep{1,2,3}{G\circ F=\Id_B}{%
    \FormulaRefAuto{FunctionExtensionality}[thm]{13,14,11}}

  \proofstep{1,2,3}{F\colon B\inj A}{%
    \FormulaRefAuto{F\circ G=\Id_B\vdash G\colon B\inj A}{1,2,15}}
\end{tabproof}

\chapter{Das Auswahlprinzip}

Band~03 setzt das Auswahlaxiom in einer reinen mengensprachlichen
Familienform an. In diesem Kapitel werden daraus die Relations- und
Surjektionsform entwickelt und der Rückweg zur Familienform gezeigt. Die
Surjektionsform ist damit keine zusätzliche Annahme, sondern die für
Funktionen besonders handliche Fassung desselben Axioms.

\section{Von der Surjektionsform zur Relationsform}

\FormulaThmDelta{%
  \pi_1\restriction_R\circ s = \Id_A\dsep x\in A
  \vdash \pi_1(s(x))=x
}{
  \DeltaRow{Mengen}{A\dsep B}
  \DeltaPrem{Totale Relationen}{\TotRel{R,A,B}}
  \DeltaPrem{Funktionen}{s\colon A\to R}
  \DeltaRow{Projektion auf erste Komponente}{\pi_1:A\times B\to A}
}
\begin{tabproofwide}
  \proofstepwidestar[1]{\pi_1\restriction_R\circ s = \Id_A}{\rA}
  \proofstepwidestar[2]{x\in A}{\rA}
  \proofstepwidestar[]{R\subseteq A\times B}{\FormulaRefAuto{F \subseteq A \times B}}
  \proofstepwidestar[2]{s(x)\in R}{\FormulaRefAuto{F\colon A\to B\dsep x\in A\vdash F(x)\in B}{2}}

  \proofstepwidestar[]{s\colon A\to R}{\rA}
  \proofstepwidestar[]{\pi_1\restriction_R\colon R\to A}{%
    \FormulaRefAuto{C\subseteq A \vdash F\restriction_C\colon C\to B}{3}}
  \proofstepwidestar[]{\pi_1\restriction_R\circ s\colon A\to A}{%
    \FormulaRefAuto{CompositionDef}[def]{5,6}}
  \proofstepwidestar[]{\Id_A\colon A\to A}{\FormulaRefAuto{IdA}[thm]}

  \proofstepwide[2]{\pi_1(s(x))}{=}{\pi_1\restriction_R(s(x))}{%
      \FormulaRefAuto{C\subseteq A \dsep x\in C \vdash F\restriction_C(x)=F(x)}{3,4}}

  \proofstepwide[2]{}{=}{(\pi_1\restriction_R\circ s)(x)}{%
      \FormulaRefAuto{x\in A \vdash (G\circ F)(x)=G(F(x))}{2}}

  \proofstepwide[1,2]{}{=}{\Id_A(x)}{%
    \FormulaRefAuto{FunctionEqualityEvaluation}[thm]{7,8,2,1}}

  \proofstepwide[2]{}{=}{x}{%
    \FormulaRefAuto{x\in A\vdash \Id_A(x)=x}{2}}

  \proofstepwide[1,2]{\pi_1(s(x))}{=}{x}{\rChain{9,10,11,12}}
\end{tabproofwide}

\FormulaThmDelta{%
  \pi_1\restriction_R\circ s = \Id_A
  \vdash \forall x\in A\,\bigl(x,\pi_2(s(x))\bigr)\in R
}{
  \DeltaRow{Mengen}{A\dsep B}
  \DeltaPrem{Totale Relationen}{\TotRel{R,A,B}}
  \DeltaPrem{Funktionen}{s\colon A\to R}
  \DeltaRow{Projektion auf erste Komponente}{\pi_1\colon A\times B\to A}
  \DeltaRow{Projektion auf zweite Komponente}{\pi_2\colon A\times B\to B}
}
\begin{tabproofwide}
  \proofstepwidestar[1]{\pi_1\restriction_R\circ s = \Id_A}{\rA}
  \proofstepwidestar[]{R\subseteq A\times B}{\FormulaRefAuto{R \subseteq A \times B}}

  \proofstepwidestar[3]{x\in A}{\rA}
  \proofstepwidestar[3]{s(x)\in R}{\FormulaRefAuto{F\colon A\to B\dsep x\in A\vdash F(x)\in B}{3}}
  \proofstepwidestar[3]{s(x)\in A\times B}{\FormulaRefAuto{A\subseteq B, x\in A\vdash x\in B}{2,4}}

  \proofstepwidestar[3]{s(x)=\bigl(\pi_1(s(x)),\pi_2(s(x))\bigr)}{%
    \FormulaRefAuto{z\in A\times B \vdash z=(\pi_1(z),\pi_2(z))}{5}}

  \proofstepwidestar[1,3]{\pi_1(s(x))=x}{%
    \FormulaRefAuto{\pi_1\restriction_R\circ s = \Id_A\dsep x\in A \vdash \pi_1(s(x))=x}{1,3}}

  \proofstepwidestar[1,3]{\bigl(\pi_1(s(x)),\pi_2(s(x))\bigr)=\bigl(x,\pi_2(s(x))\bigr)}{%
    \rIE{7,6}}

  \proofstepwidestar[1,3]{s(x)=\bigl(x,\pi_2(s(x))\bigr)}{%
    \FormulaRefAuto{a=b, b=c\vdash a=c}{6,8}}

  \proofstepwidestar[1,3]{\bigl(x,\pi_2(s(x))\bigr)\in R}{%
    \rIE{9,4}}

  \proofstepwidestar[1]{\forall x\in A\,\bigl(x,\pi_2(s(x))\bigr)\in R}{%
    \rUI{\rRI{3,10}}}
\end{tabproofwide}


\FormulaThmDelta{%
  \pi_1\restriction_R\circ s = \Id_A
  \vdash \forall x\in A\,\bigl(x,\pi_2\restriction_R(s(x))\bigr)\in R
}{
  \DeltaRow{Mengen}{A\dsep B}
  \DeltaPrem{Totale Relationen}{\TotRel{R,A,B}}
  \DeltaPrem{Funktionen}{s\colon A\to R}
  \DeltaRow{Projektion auf zweite Komponente}{\pi_2\colon A\times B\to B}
}
\begin{tabproof}
  \proofstep{1}{\pi_1\restriction_R\circ s = \Id_A}{\rA}
  \proofstep{}{R\subseteq A\times B}{\FormulaRefAuto{R \subseteq A \times B}}

  \proofstep{1}{\forall x\in A\,\bigl(x,\pi_2(s(x))\bigr)\in R}{%
    \FormulaRefAuto{\pi_1\restriction_R\circ G = \Id_A \vdash \forall x\in A\,\bigl(x,\pi_2(G(x))\bigr)\in R}{1}}

  \proofstep{}{x\in A}{\rA}
  \proofstep{}{s(x)\in R}{\FormulaRefAuto{F\colon A\to B\dsep x\in A\vdash F(x)\in B}{4}}

  \proofstep{}{ \pi_2\restriction_R(s(x))=\pi_2(s(x))}{%
    \FormulaRefAuto{C\subseteq A \dsep x\in C \vdash F\restriction_C(x)=F(x)}{2,5}}

  \proofstep{}{ \pi_2(s(x))=\pi_2\restriction_R(s(x))}{%
    \FormulaRefAuto{a=b\vdash b=a}{6}}

  \proofstep{}{(x,\pi_2(s(x)))\in R}{\rUE{3}}
  \proofstep{}{(x,\pi_2\restriction_R(s(x)))\in R}{\rIE{7,8}}

  \proofstep{1}{\forall x\in A\,\bigl(x,\pi_2\restriction_R(s(x))\bigr)\in R}{%
    \rUI{\rRI{4,9}}}
\end{tabproof}

\FormulaThmDelta{%
  \pi_1\restriction_R\circ s = \Id_A
  \vdash \exists F\colon A\to B\,\forall x\in A\,(x,F(x))\in R
}{
  \DeltaRow{Mengen}{A\dsep B}
  \DeltaPrem{Totale Relationen}{\TotRel{R,A,B}}
  \DeltaPrem{Funktionen}{s\colon A\to R}
  \DeltaRow{Projektion auf zweite Komponente}{\pi_2\colon A\times B\to B}
}
\begin{tabproof}
  \proofstep{1}{\pi_1\restriction_R\circ s = \Id_A}{\rA}

  \proofstep{2}{\forall x\in A\,\bigl(x,\pi_2\restriction_R(s(x))\bigr)\in R}{%
    \FormulaRefAuto{\pi_1\restriction_R\circ G = \Id_A \vdash \forall x\in A\,\bigl(x,\pi_2\restriction_R(G(x))\bigr)\in R}{1}}

  \proofstep{3}{\pi_2\restriction_R\circ s\colon A\to B}{%
    \FormulaRefAuto{G\circ F\colon A\to C}}

  \proofstep{4}{x\in A}{\rA}

  \proofstep{4}{\bigl(x,\pi_2\restriction_R(s(x))\bigr)\in R}{\rUE{2}}

  \proofstep{4}{(\pi_2\restriction_R\circ s)(x)=\pi_2\restriction_R(s(x))}{%
    \FormulaRefAuto{x\in A \vdash (G\circ F)(x)=G(F(x))}{4}}

  \proofstep{4}{\pi_2\restriction_R(s(x))=(\pi_2\restriction_R\circ s)(x)}{%
    \FormulaRefAuto{a=b\vdash b=a}{6}}

  \proofstep{4}{\bigl(x,(\pi_2\restriction_R\circ s)(x)\bigr)\in R}{\rIE{7,5}}

  \proofstep{4}{\forall x\in A\,\bigl(x,(\pi_2\restriction_R\circ s)(x)\bigr)\in R}{%
    \rUI{\rRI{4,8}}}

  \proofstep{1}{\exists F\colon A\to B\,\forall x\in A\,(x,F(x))\in R}{%
    \rEI{\rAI{9,3}}}
\end{tabproof}

\FormulaThmDelta{%
  \exists s\colon A\to R\,\bigl(\pi_1\restriction_R\circ s = \Id_A\bigr)
  \vdash \exists F\colon A\to B\,\forall x\in A\,\bigl((x,F(x))\in R\bigr)
}{
  \DeltaRow{Mengen}{A\dsep B}
  \DeltaPrem{Totale Relationen}{\TotRel{R,A,B}}
  \DeltaPrem{Funktionen}{s\colon A\to R}
  \DeltaRow{Projektion auf zweite Komponente}{\pi_2\colon A\times B\to B}
}
\begin{tabproof}
  \proofstep{1}{\exists s\colon A\to R\,(\pi_1\restriction_R\circ s = \Id_A)}{\rA}
  \proofstep{2}{s\colon A\to R\land \pi_1\restriction_R\circ s = \Id_A}{\rA}

  \proofstep{2}{\exists F\colon A\to B\,\forall x\in A\,((x,F(x))\in R)}{%
    \FormulaRefAuto{\pi_1\restriction_R\circ s = \Id_A \vdash \exists F\colon A\to B\,\forall x\in A\,(x,F(x))\in R}{2}}

  \proofstep{1}{\exists F\colon A\to B\,\forall x\in A\,((x,F(x))\in R)}{%
    \rEE{1,2,3}}
\end{tabproof}

\section{Von der Relationsform zur Familienform}

\FormulaThmDelta{%
  \forall x\in M\,\bigl((x,F(x))\in \MemRel(M)\bigr)\dsep X\in M\vdash F(X)\in X
}{
  \DeltaRow{Mengen}{M \dsep X \dsep a}
  \DeltaPrem{\makecell[l]{Paarweise disjunkte\\ nichtleere Familie}}{\DisjFam(M)}
  \DeltaPrem{Funktionen}{F\colon M\to \bigcup M}
  \DeltaRow{Mitgliedschaftsrelation}{\MemRel(M)}[
    \FormulaRefAuto{\DisjFam(M)\vdash \TotRel{\MemRel(M),M,\bigcup M}}]
}
\begin{tabproof}
  \proofstep{1}{%
    \forall x\in M\,\bigl((x,F(x))\in \MemRel(M)\bigr)}{%
    \rA}

  \proofstep{2}{%
    X\in M}{%
    \rA}

  \proofstep{1,2}{%
    (X,F(X))\in \MemRel(M)}{%
    \rRE{2,\rUE{1}}}

  \proofstep{1,2}{%
    X\in M \land F(X)\in \bigcup M \land F(X)\in X}{%
    \FormulaRefAuto{(X,a)\in \MemRel(M) \vdash X\in M \land a\in \bigcup M \land a\in X}{3}}

  \proofstep{1,2}{%
    F(X)\in X}{%
    \rAEb{4}}
\end{tabproof}

\FormulaThmDelta[Existenz eines Auswahlelements in $X$]{%
  \forall x\in M\,\bigl((x,F(x))\in \MemRel(M)\bigr)\dsep X\in M
  \vdash \exists a\in X\,\bigl(a\in F[M]\bigr)
}{
  \DeltaRow{Mengen}{M \dsep X}
  \DeltaPrem{\makecell[l]{Paarweise disjunkte\\ nichtleere Familie}}{\DisjFam(M)}
  \DeltaPrem{Funktionen}{F\colon M\to \bigcup M}
}
\begin{tabproof}
  \proofstep{1}{F\colon M\to \bigcup M}{\rA}
  \proofstep{2}{\forall x\in M\,\bigl((x,F(x))\in \MemRel(M)\bigr)}{\rA}
  \proofstep{3}{X\in M}{\rA}

  \proofstep{2,3}{F(X)\in X}{%
    \FormulaRefAuto{\forall x\in M\,\bigl((x,F(x))\in \MemRel(M)\bigr)\dsep X\in M\vdash F(X)\in X}{2,3}}

  \proofstep{}{M\subseteq M}{\FormulaRefAuto{A\subseteq A}}
  \proofstep{1,3}{F(X)\in F[M]}{%
    \FormulaRefAuto{C\subseteq A\dsep F\colon A\to B\dsep x\in C\vdash F(x)\in F[C]}{5,1,3}}

  \proofstep{1,2,3}{\exists a\in X\,\bigl(a\in F[M]\bigr)}{\rEI{\rAI{4,6}}}
\end{tabproof}

\FormulaThmDelta[Auswahlelement in $X$ ist $F(X)$]{%
  \forall x\in M\,\bigl((x,F(x))\in \MemRel(M)\bigr)\dsep X\in M\dsep a\in X\dsep a\in F[M]
  \vdash a=F(X)
}{
  \DeltaRow{Mengen}{M \dsep X \dsep a \dsep Y}
  \DeltaPrem{\makecell[l]{Paarweise disjunkte\\ nichtleere Familie}}{\DisjFam(M)}
  \DeltaPrem{Funktionen}{F\colon M\to \bigcup M}
  \DeltaRow{Mitgliedschaftsrelation}{\MemRel(M)}[
    \FormulaRefAuto{\DisjFam(M)\vdash \TotRel{\MemRel(M),M,\bigcup M}}]
}
\begin{tabproof}
  \proofstep{1}{F\colon M\to \bigcup M}{\rA}
  \proofstep{2}{\forall x\in M\,\bigl((x,F(x))\in \MemRel(M)\bigr)}{\rA}
  \proofstep{3}{X\in M}{\rA}
  \proofstep{4}{a\in X}{\rA}
  \proofstep{5}{a\in F[M]}{\rA}

  \proofstep{}{M\subseteq M}{\FormulaRefAuto{A\subseteq A}}
  \proofstep{1,5}{\exists Y\in M\,a=F(Y)}{%
    \FormulaRefAuto{C\subseteq A\dsep F\colon A\to B\dsep y\in F[C]\vdash \exists x\in C\,y=F(x)}{6,1,5}}

  \proofstep{8}{Y\in M\land a=F(Y)}{\rA}
  \proofstep{8}{Y\in M}{\rAEa{8}}
  \proofstep{8}{a=F(Y)}{\rAEb{8}}

  \proofstep{2,9}{F(Y)\in Y}{%
    \FormulaRefAuto{\forall x\in M\,\bigl((x,F(x))\in \MemRel(M)\bigr)\dsep X\in M\vdash F(X)\in X}{2,9}}

  \proofstep{8}{a\in Y}{%
    \FormulaRefAuto{a=b \dsep b\in A \vdash a\in A}{10,11}}

  \proofstep{3,4,9,12}{X=Y}{%
    \FormulaRefAuto{X \in M \dsep Y \in M \dsep a \in X \dsep a \in Y \vdash X = Y}{3,9,4,12}}

  \proofstep{8}{a=F(X)}{%
    \FormulaRefAuto{x = y \dsep a = F(y) \vdash a = F(x)}{13,10}}

  \proofstep{1,2,3,4,5}{a=F(X)}{\rEE{7,8,14}}
\end{tabproof}


\FormulaThmDelta[Eindeutige Existenz eines Auswahlelements in $X$]{%
  \forall x\in M\,\bigl((x,F(x))\in \MemRel(M)\bigr)\dsep X\in M
  \vdash \exists! a\in X\,\bigl(a\in F[M]\bigr)
}{
  \DeltaRow{Mengen}{M \dsep X}
  \DeltaPrem{\makecell[l]{Paarweise disjunkte\\ nichtleere Familie}}{\DisjFam(M)}
  \DeltaPrem{Funktionen}{F\colon M\to \bigcup M}
  \DeltaRow{Mitgliedschaftsrelation}{\MemRel(M)}[
    \FormulaRefAuto{\DisjFam(M)\vdash \TotRel{\MemRel(M),M,\bigcup M}}]
}
\begin{tabproofsplit}
\proofpart{Existenz}
\proofstep{1}{\forall x\in M\,\bigl((x,F(x))\in \MemRel(M)\bigr)}{\rA}
  \proofstep{2}{X\in M}{\rA}

  \proofstep{1,2}{\exists a\in X\,\bigl(a\in F[M]\bigr)}{%
    \FormulaRefAuto{\forall x\in M\,\bigl((x,F(x))\in \MemRel(M)\bigr)\dsep X\in M
    \vdash \exists a\in X\,\bigl(a\in F[M]\bigr)}{1,2}}


\closeproofpart

\proofpart{Eindeutigkeit}
  \setcounter{proofstepnr}{3}% Fortlaufende Schritte dieses Beweises.
\proofstep{4}{b\in X\land b\in F[M]}{\rA}
  \proofstep{5}{c\in X\land c\in F[M]}{\rA}

  \proofstep{4}{b\in X}{\rAEa{4}}
  \proofstep{4}{b\in F[M]}{\rAEb{4}}
  \proofstep{1,2,4}{b=F(X)}{%
    \FormulaRefAuto{\forall x\in M\,\bigl((x,F(x))\in \MemRel(M)\bigr)\dsep X\in M\dsep a\in X\dsep a\in F[M]
    \vdash a=F(X)}{1,2,6,7}}

  \proofstep{5}{c\in X}{\rAEa{5}}
  \proofstep{5}{c\in F[M]}{\rAEb{5}}
  \proofstep{1,2,5}{c=F(X)}{%
    \FormulaRefAuto{\forall x\in M\,\bigl((x,F(x))\in \MemRel(M)\bigr)\dsep X\in M\dsep a\in X\dsep a\in F[M]
    \vdash a=F(X)}{1,2,9,10}}

  \proofstep{1,2,4,5}{b=c}{\FormulaRefAuto{a=b,c=b\vdash a=c}{8,11}}


\closeproofpart

\proofpart{Eindeutige Existenz}
  \setcounter{proofstepnr}{12}% Fortlaufende Schritte dieses Beweises.
\proofstep{1,2}{\exists! a\in X\,\bigl(a\in F[M]\bigr)}{\UEI{3,4,5,12}}

\closeproofpart
\end{tabproofsplit}


\FormulaThmDelta[Auswahlfunktion liefert Auswahlmenge]{%
  \exists F\colon M\to \bigcup M\,\forall x\in M\,\bigl((x,F(x))\in \MemRel(M)\bigr)
  \vdash \exists L\,\forall X\in M\,\exists! a\in X\,\bigl(a\in L\bigr)
}{
  \DeltaRow{Mengen}{M}
  \DeltaPrem{\makecell[l]{Paarweise disjunkte\\ nichtleere Familie}}{\DisjFam(M)}
  \DeltaRow{Mitgliedschaftsrelation}{\MemRel(M)}[
    \FormulaRefAuto{\DisjFam(M)\vdash \TotRel{\MemRel(M),M,\bigcup M}}]
}
\begin{tabproof}
  \proofstep{1}{%
    \exists F\colon M\to \bigcup M\,\forall x\in M\,\bigl((x,F(x))\in \MemRel(M)\bigr)
  }{\rA}

  \proofstep{2}{%
    F\colon M\to \bigcup M \land \forall x\in M\,\bigl((x,F(x))\in \MemRel(M)\bigr)
  }{\rA}

  \proofstep{2}{%
    \forall x\in M\,\bigl((x,F(x))\in \MemRel(M)\bigr)
  }{\rAEb{2}}

  \proofstep{4}{X\in M}{\rA}

  \proofstep{2,4}{%
    \exists! a\in X\,\bigl(a\in F[M]\bigr)
  }{%
    \FormulaRefAuto{\forall x\in M\,\bigl((x,F(x))\in \MemRel(M)\bigr)\dsep X\in M
    \vdash \exists! a\in X\,\bigl(a\in F[M]\bigr)}{3,4}
  }

  \proofstep{2}{%
    \forall X\in M\,\exists! a\in X\,\bigl(a\in F[M]\bigr)
  }{\rUI{\rRI{4,5}}}

  \proofstep{2}{%
    \exists L\,\forall X\in M\,\exists! a\in X\,\bigl(a\in L\bigr)
  }{\rEI{6}}

  \proofstep{1}{%
    \exists L\,\forall X\in M\,\exists! a\in X\,\bigl(a\in L\bigr)
  }{\rEE{1,2,7}}
\end{tabproof}



\section{Von der Familienform zur Surjektionsform}

\FormulaThmDeltaR{%
  \begin{aligned}[t]
    &\forall M\,\Bigl(\DisjFam(M)\rightarrow
      \exists L\,\forall X\in M\,\exists! a\in X\,(a\in L)\Bigr)\\
    &\vdash\ \exists L\,\forall X\in \Fib_G[A]\,\exists! a\,(a\in X \land a\in L)
  \end{aligned}%
}{%
  \forall M\,\bigl(\DisjFam(M)\rightarrow \exists L\,\forall X\in M\,\exists! a\in X\,a\in L\bigr)
  \vdash \exists L\,\forall X\in \Fib_G[A]\,\exists! a\,(a\in X \land a\in L)
}{%
  \DeltaDecl{Mengen}{A\dsep B\dsep X\dsep a}%
  \DeltaPrem{Surjektive Funktion}{G\colon B\sur A}%
}
\begin{tabproof}
  \proofstep{1}{\forall M\,\bigl(\DisjFam(M)\rightarrow \exists L\,\forall X\in M\,\exists! a\in X\,(a\in L)\bigr)}{\rA}
  \proofstep{}{\DisjFam(\Fib_G[A])}{\FormulaRefAuto{\DisjFam(\Fib_F[B])}}
  \proofstep{1}{\exists L\,\forall X\in \Fib_G[A]\,\exists! a\,(a\in X \land a\in L)}{\rRE{\rUE{1},2}}
\end{tabproof}

\FormulaThmDeltaR{%
  \begin{aligned}[t]
    &\forall X\in \Fib_G[A]\,\exists! a\,(a\in X \land a\in L)\dsep x\in A\\
    &\vdash\ \exists! a\,\bigl(a\in \Fib_G(x)\land a\in L\bigr)
  \end{aligned}%
}{%
  \forall X\in \Fib_G[A]\,\exists! a\,(a\in X \land a\in L)\dsep x\in A
  \vdash \exists! a\,\bigl(a\in \Fib_G(x)\land a\in L\bigr)%
}{%
  \DeltaRow{Mengen}{A\dsep B\dsep L\dsep x}%
  \DeltaPrem{Funktionen}{G\colon B\to A}%
}
\begin{tabproof}
  \proofstep{1}{G\colon B\to A}{\rA}
  \proofstep{2}{\forall X\in \Fib_G[A]\,\exists! a\,(a\in X \land a\in L)}{\rA}
  \proofstep{3}{x\in A}{\rA}
  \proofstep{1}{\Fib_G\colon A\to \powerset(B)}{%
    \FormulaRefAuto{\Fib_G\colon A\to \powerset(B)}{1}}
  \proofstep{}{A\subseteq A}{\FormulaRefAuto{A\subseteq A}}
  \proofstep{1,3}{\Fib_G(x)\in \Fib_G[A]}{%
    \FormulaRefAuto{C\subseteq A\dsep F\colon A\to B\dsep x\in C\vdash F(x)\in F[C]}{5,4,3}}
  \proofstep{1,2,3}{\exists! a\,\bigl(a\in \Fib_G(x)\land a\in L\bigr)}{%
    \rRE{\rUE{2},6}}
\end{tabproof}



% ============================================================
% 3.15.11.x  Konstruktion einer Sektion aus einer Auswahlmenge
% (Einfügen z.B. nach dem Lemma
%  \forall X\in \Fib_G[A]\,\exists! a\,(a\in X \land a\in L) \dsep x\in A
%  \vdash \exists! b\,(b\in \Fib_G(x)\land b\in L)
% und vor dem Hauptsatz \ACsur.)
% ============================================================

\subsection{Die Sektion zu einer surjektiven Funktion}

\subsubsection{Direkte Definition der Sektion}

\FormulaDefDeltaK[Sektion zu einer Auswahlmenge]{%
  \forall X\in \Fib_G[A]\,\exists! a\,(a\in X\land a\in L)
  \vdash
  \begin{aligned}[t]
    &\Sec_{G,L}\colon A\to B,\\[-2pt]
    &\forall x\in A\,
      \Bigl(\Sec_{G,L}(x)\coloneqq
        \iota b\bigl(b\in\Fib_G(x)\land b\in L\bigr)\Bigr)
  \end{aligned}%
}{ChoiceSectionDef}{%
  \DeltaRow{Mengen}{A\dsep B\dsep L\dsep X\dsep a\dsep b\dsep x}
  \DeltaRow{Surjektive Funktion}{G\colon B\sur A}
  \DeltaRow{Funktionen}{\Sec_{G,L}}
}
\begin{tabproof}
  \proofstep{1}{G\colon B\sur A}{\rA}
  \proofstep{1}{G\colon B\to A}{%
    \FormulaRefAuto{Surjektive Funktion}[def]{1}}
  \proofstep{3}{%
    \forall X\in\Fib_G[A]\,\exists! a\,(a\in X\land a\in L)}{\rA}
  \proofstep{4}{x\in A}{\rA}
  \proofstep{1,3,4}{%
    \begin{aligned}[t]
      &\iota b\bigl(b\in\Fib_G(x)\land b\in L\bigr)\in\Fib_G(x)\land{}\\[-2pt]
      &\iota b\bigl(b\in\Fib_G(x)\land b\in L\bigr)\in L
    \end{aligned}%
  }{%
    \FormulaRefAuto{%
      \forall X\in\Fib_G[A]\,\exists! a\,(a\in X\land a\in L)
      \dsep x\in A
      \vdash \exists! a\,(a\in\Fib_G(x)\land a\in L)%
    }[thm]{2,3,4}}
  \proofstep{1,3,4}{%
    \iota b\bigl(b\in\Fib_G(x)\land b\in L\bigr)\in\Fib_G(x)}{%
    \rAEa{5}}
  \proofstep{1,3,4}{%
    \iota b\bigl(b\in\Fib_G(x)\land b\in L\bigr)\in B}{%
    \FormulaRefAuto{y\in B\dsep x\in\Fib_F(y)\vdash x\in A}[thm]{2,4,6}}
\end{tabproof}

Der innere \(\iota\)-Term bezeichnet gemäß der Iota-Konvention die durch
die Auswahlvoraussetzung eindeutig bestimmte Auswahl aus der Faser; nur die
äußere Funktions-Iota-Konstruktion entfällt.


\FormulaThmDelta[Sektion ist eine Funktion]{%
  \forall X\in \Fib_G[A]\,\exists! a\,(a\in X \land a\in L)
  \vdash \Sec_{G,L}\colon A\to B%
}{
  \DeltaRow{Mengen}{A\dsep B\dsep L\dsep X\dsep a}
  \DeltaRow{Surjektive Funktion}{G\colon B\sur A}
}
\begin{tabproof}
  \proofstep{1}{\forall X\in \Fib_G[A]\,\exists! a\,(a\in X \land a\in L)}{\rA}
  \proofstep{1}{\Sec_{G,L}\colon A\to B}{%
    \FormulaRefAuto{ChoiceSectionDef}[def]{1}}
\end{tabproof}

\FormulaThmDelta[Sektionsgleichung]{%
  \forall X\in \Fib_G[A]\,\exists! a\,(a\in X \land a\in L) \dsep x\in A
  \vdash G\bigl(\Sec_{G,L}(x)\bigr)=x%
}{
  \DeltaRow{Mengen}{A\dsep B\dsep L\dsep X\dsep a\dsep x}
  \DeltaRow{Surjektive Funktion}{G\colon B\sur A}
  \DeltaRow{Funktionen}{\Sec_{G,L}}
}
\begin{tabproof}
  \proofstep{1}{\forall X\in \Fib_G[A]\,\exists! a\,(a\in X \land a\in L)}{\rA}
  \proofstep{2}{x\in A}{\rA}
  \proofstep{}{G\colon B\to A}{%
    \FormulaRefAuto{Surjektive Funktion}[def]}
  \proofstep{1}{\forall u\in A\,
    \Bigl(\Sec_{G,L}(u)=
      \iota b\bigl(b\in\Fib_G(u)\land b\in L\bigr)\Bigr)}{%
    \FormulaRefAuto{ChoiceSectionDef}[def]{1}}
  \proofstep{1,2}{\Sec_{G,L}(x)=
    \iota b\bigl(b\in\Fib_G(x)\land b\in L\bigr)}{%
    \rRE{\rUE{4},2}}
  \proofstep{1,2}{%
    \begin{aligned}[t]
      &\iota b\bigl(b\in\Fib_G(x)\land b\in L\bigr)
        \in\Fib_G(x)\land{}\\[-2pt]
      &\iota b\bigl(b\in\Fib_G(x)\land b\in L\bigr)\in L
    \end{aligned}%
  }{%
    \FormulaRefAuto{%
      \forall X\in\Fib_G[A]\,\exists! a\,(a\in X\land a\in L)
      \dsep x\in A
      \vdash \exists! a\,(a\in\Fib_G(x)\land a\in L)%
    }[thm]{3,1,2}}
  \proofstep{1,2}{%
    \iota b\bigl(b\in\Fib_G(x)\land b\in L\bigr)\in\Fib_G(x)}{%
    \rAEa{6}}
  \proofstep{1,2}{\Sec_{G,L}(x)\in \Fib_G(x)}{%
    \rIE{5,7}}
  \proofstep{1,2}{G\bigl(\Sec_{G,L}(x)\bigr)=x}{%
    \FormulaRefAuto{x\in A \dsep b\in \Fib_G(x) \vdash G(b)=x}{2,8}}
\end{tabproof}

\FormulaThmDeltaK[Auswahlmenge liefert Rechtsinverse]{%
  \forall X\in \Fib_G[A]\,\exists! a\,(a\in X \land a\in L)
  \vdash \exists F\colon A\to B\,\bigl(G\circ F=\Id_A\bigr)%
}{ChoiceSetRightInverse}{
  \DeltaRow{Mengen}{A\dsep B\dsep L\dsep X\dsep a\dsep x}
  \DeltaRow{Surjektive Funktion}{G\colon B\sur A}
}
\begin{tabproof}
  \proofstep{1}{\forall X\in \Fib_G[A]\,\exists! a\,(a\in X \land a\in L)}{\rA}

  \proofstep{1}{\Sec_{G,L}\colon A\to B}{%
    \FormulaRefAuto{%
      \forall X\in \Fib_G[A]\,\exists! a\,(a\in X \land a\in L)
      \vdash \Sec_{G,L}\colon A\to B}{1}}

  \proofstep{3}{x\in A}{\rA}

  \proofstep{1,3}{\bigl(G\circ \Sec_{G,L}\bigr)(x)=G\bigl(\Sec_{G,L}(x)\bigr)}{%
    \FormulaRefAuto{x\in A \vdash (G\circ F)(x)=G(F(x))}{2,3}}

  \proofstep{1,3}{G\bigl(\Sec_{G,L}(x)\bigr)=x}{%
    \FormulaRefAuto{%
      \forall X\in \Fib_G[A]\,\exists! a\,(a\in X \land a\in L) \dsep x\in A
      \vdash G\bigl(\Sec_{G,L}(x)\bigr)=x}{1,3}}

  \proofstep{1,3}{\bigl(G\circ \Sec_{G,L}\bigr)(x)=x}{%
    \FormulaRefAuto{a=b,\, b=c \vdash a=c}{4,5}}

  \proofstep{3}{\Id_A(x)=x}{\FormulaRefAuto{x\in A\vdash \Id_A(x)=x}{3}}
  \proofstep{3}{x=\Id_A(x)}{\FormulaRefAuto{a=b\vdash b=a}{7}}

  \proofstep{1,3}{\bigl(G\circ \Sec_{G,L}\bigr)(x)=\Id_A(x)}{%
    \FormulaRefAuto{a=b,\, b=c \vdash a=c}{6,8}}

  \proofstep{1}{x\in A\rightarrow \bigl(G\circ \Sec_{G,L}\bigr)(x)=\Id_A(x)}{\rRI{3,9}}
  \proofstep{1}{\forall x\in A\,\bigl(G\circ \Sec_{G,L}\bigr)(x)=\Id_A(x)}{\rUI{10}}

  \proofstep{}{G\colon B\to A}{\FormulaRefAuto{Surjektive Funktion}[def]}
  \proofstep{1}{G\circ \Sec_{G,L}\colon A\to A}{%
    \FormulaRefAuto{CompositionDef}[def]{2,12}}
  \proofstep{}{\Id_A\colon A\to A}{\FormulaRefAuto{IdA}[thm]}
  \proofstep{1}{G\circ \Sec_{G,L}=\Id_A}{%
    \FormulaRefAuto{FunctionExtensionality}[thm]{13,14,11}}

  \proofstep{1}{\Sec_{G,L}\colon A\to B \land G\circ \Sec_{G,L}=\Id_A}{\rAI{2,15}}
  \proofstep{1}{\exists F\colon A\to B\,\bigl(G\circ F=\Id_A\bigr)}{\rEI{16}}
\end{tabproof}


% -------------------------------------------------
% 10) Hauptsatz: ACfam ⇒ ACsur (für eine gegebene Surjektion G)
% -------------------------------------------------
\FormulaThmDeltaR{%
  \begin{aligned}[t]
    &\forall M\,\Bigl(\DisjFam(M)\rightarrow
      \exists L\,\forall X\in M\,\exists! a\,(a\in X \land a\in L)\Bigr)\\
    &\vdash\ \exists F\colon A\to B\,\bigl(G\circ F=\Id_A\bigr)
  \end{aligned}%
}{%
  \forall M\,\bigl(\DisjFam(M)\rightarrow \exists L\,\forall X\in M\,\exists! a\,(a\in X \land a\in L)\bigr)
  \vdash \exists F\colon A\to B\,\bigl(G\circ F=\Id_A\bigr)%
}{%
  \DeltaRow{Mengen}{A\dsep B\dsep x}%
  \DeltaRow{Surjektive Funktion}{G\colon B\sur A}%
}
\begin{tabproof}
  \proofstep{1}{\forall M\,\bigl(\DisjFam(M)\rightarrow \exists L\,\forall X\in M\,\exists! a\,(a\in X \land a\in L)\bigr)}{\rA}

  \proofstep{1}{\exists L\,\forall X\in \Fib_G[A]\,\exists! a\,(a\in X \land a\in L)}{%
    \FormulaRefAuto{%
      \forall M\,\bigl(\DisjFam(M)\rightarrow \exists L\,\forall X\in M\,\exists! a\in X\, a\in L\bigr)
      \vdash \exists L\,\forall X\in \Fib_G[A]\,\exists! a\,(a\in X \land a\in L)
    }{1}}

  \proofstep{3}{\forall X\in \Fib_G[A]\,\exists! a\,(a\in X \land a\in L)}{\rA}
  \proofstep{3}{\exists F\colon A\to B\,\bigl(G\circ F=\Id_A\bigr)}{%
    \FormulaRefAuto{%
      \forall X\in \Fib_G[A]\,\exists! a\,(a\in X \land a\in L)
      \vdash \exists F\colon A\to B\,\bigl(G\circ F=\Id_A\bigr)
    }{3}}
  \proofstep{1}{\exists F\colon A\to B\,\bigl(G\circ F=\Id_A\bigr)}{\rEE{2,3,4}}
\end{tabproof}

\subsection{Äquivalente axiomatische Funktionsfassung}

Die vorangehende Konstruktion gewinnt aus der funktionsfreien Familienform
für jede Surjektion eine Rechtsinverse. Deshalb kann das Auswahlaxiom in
diesem Band äquivalent wie folgt registriert werden.

\FormulaAxiomDeltaK[Auswahlprinzip (AC -- Surjektionsform)]{\exists F\colon A \to B\,G\circ F = \Id_A}{Auswahlaxiom}{
  \DeltaDecl{Mengen}{A \dsep B}
  \DeltaPrem{Surjektive Funktion}{G\colon B \sur A}
}

\FormulaDefDeltaK[Auswahlprinzip (globale Surjektionsform)]{%
  \ACsur\coloneqq
  \forall A\,\forall B\,\forall G\,
  \bigl(G\colon B\sur A\rightarrow
  \exists F\colon A\to B\,G\circ F=\Id_A\bigr)%
}{ACsur}{
  \DeltaDecl{Mengen}{A \dsep B}
  \DeltaDecl{Funktionen}{F \dsep G}
  \DeltaPrem{Surjektive Funktion}{G\colon B\sur A}
}

Die Rückrichtung wird im folgenden Abschnitt über die Mitgliedschaftsrelation
einer paarweise disjunkten nichtleeren Familie hergeleitet. Zusammen lauten
die beiden Richtungen kompakt
\[
  \forall M\,\Bigl(
    \DisjFam(M)\rightarrow
    \exists L\,\forall X\in M\,\exists!a\,
    (a\in X\land a\in L)
  \Bigr)
  \quad\Longleftrightarrow\quad
  \ACsur .
\]

\subsection{Folgerungen}

\FormulaThmDelta[AC -- Relationsform aus der Surjektionsform]{%
  \exists F\colon A \to B\,\forall x\in A\,\bigl((x,F(x))\in R\bigr)%
}{
  \DeltaDecl{Mengen}{A \dsep B}
  \DeltaPrem{Totale Relation}{\TotRel{R,A,B}}
  \DeltaDecl{Relation}{R}
}
\begin{tabproof}
  \proofstep{}{%
    \pi_1\restriction_R\colon R\sur A%
  }{%
    \FormulaRefAuto{\pi_1\restriction_R\colon R\sur A}%
  }

  \proofstep{1}{%
    \exists s\colon A\to R\,\bigl(\pi_1\restriction_R\circ s=\Id_A\bigr)%
  }{%
    \FormulaRefAuto{\exists F\colon A \to B\,G\circ F = \Id_A}{1}%
  }

  \proofstep{2}{%
    \exists F\colon A \to B\,\forall x\in A\,\bigl((x,F(x))\in R\bigr)%
  }{%
    \FormulaRefAuto{\exists s\colon A\to R\,\bigl(\pi_1\restriction_R\circ s=\Id_A\bigr)\vdash
    \exists F\colon A\to B\,\forall x\in A\,\bigl((x,F(x))\in R\bigr)}{2}%
  }
\end{tabproof}

\FormulaThmDelta[AC -- Familienform aus der Surjektionsform]{%
  \exists L\,\forall X\in M\,\exists! a\in X\,\bigl(a\in L\bigr)%
}{
  \DeltaDecl{Mengen}{M \dsep X \dsep a \dsep L}
  \DeltaPrem{\makecell[l]{Paarweise disjunkte\\ nichtleere Familie}}{\DisjFam(M)}
}
\begin{tabproof}
  \proofstep{1}{\DisjFam(M)}{\rA}

  \proofstep{1}{%
    \TotRel{\MemRel(M),M,\bigcup M}%
  }{%
    \FormulaRefAuto{%
      \DisjFam(M)\vdash \TotRel{\MemRel(M),M,\bigcup M}%
    }{1}%
  }

  \proofstep{1}{%
    \exists F\colon M\to \bigcup M\,\forall x\in M\,\bigl((x,F(x))\in \MemRel(M)\bigr)%
  }{%
    \FormulaRefAuto{\exists F\colon A \to B\,\forall x\in A\,\bigl((x,F(x))\in R\bigr)}{2}%
  }

  \proofstep{1}{%
    \exists L\,\forall X\in M\,\exists! a\in X\,\bigl(a\in L\bigr)%
  }{%
    \FormulaRefAuto{\exists F\colon M\to \bigcup M\,\forall x\in M\,\bigl((x,F(x))\in \MemRel(M)\bigr)
    \vdash \exists L\,\forall X\in M\,\exists! a\in X\,\bigl(a\in L\bigr)}{3}%
  }
\end{tabproof}


\FormulaThmDelta[Rechtsinverse surjektiver Funktionen]{%
  F\colon A\sur B\vdash
  \exists H\colon B\inj A\,\forall y\in B\,\bigl(F(H(y))=y\bigr)
}{%
  \DeltaDecl{Mengen}{A\dsep B\dsep y}%
  \DeltaDecl{Funktionen}{F\dsep H}%
}
\begin{tabproofwide}
  \proofstepwidestar[1]{F\colon A\sur B}{\rA}
  \proofstepwidestar[1]{F\colon A\to B}{\FormulaRefAuto{Surjektive Funktion}{1}}

  \proofstepwidestar[1]{\exists H\colon B\to A\,F\circ H=\Id_B}{%
    \FormulaRefAuto{Auswahlaxiom}{1}}

  \proofstepwidestar[4]{H\colon B\to A \land F\circ H=\Id_B}{\rA}
  \proofstepwidestar[4]{H\colon B\to A}{\rAEa{4}}
  \proofstepwidestar[4]{F\circ H=\Id_B}{\rAEb{4}}

  \proofstepwidestar[1,4]{H\colon B\inj A}{%
    \FormulaRefAuto{F\circ G=\Id_B\vdash G\colon B\inj A}{2,5,6}}

  \proofstepwidestar[8]{y\in B}{\rA}

  \proofstepwidestar[1,4]{F\circ H\colon B\to B}{%
    \FormulaRefAuto{CompositionDef}[def]{5,2}}
  \proofstepwidestar[]{\Id_B\colon B\to B}{\FormulaRefAuto{IdA}[thm]}

  \proofstepwide[1,4,8]{F(H(y))}{=}{(F\circ H)(y)}{%
    \FormulaRefAuto{x\in A \vdash F(G(x))=(F\circ G)(x)}{8}}

  \proofstepwide[4,8]{}{=}{\Id_B(y)}{%
    \FormulaRefAuto{FunctionEqualityEvaluation}[thm]{9,10,8,6}}

  \proofstepwide[8]{}{=}{y}{%
    \FormulaRefAuto{x\in A\vdash \Id_A(x)=x}{8}}

  \proofstepwide[1,4,8]{F(H(y))}{=}{y}{\rChain{11,12,13}}

  \proofstepwidestar[1,4]{\forall y\in B\,\bigl(F(H(y))=y\bigr)}{%
    \rUI{\rRI{8,14}}}

  \proofstepwidestar[1,4]{%
    H\colon B\inj A \land \forall y\in B\,\bigl(F(H(y))=y\bigr)
  }{\rAI{7,15}}

  \proofstepwidestar[1,4]{%
    \exists H\colon B\inj A\,\forall y\in B\,\bigl(F(H(y))=y\bigr)
  }{\rEI{16}}

  \proofstepwidestar[1]{%
    \exists H\colon B\inj A\,\forall y\in B\,\bigl(F(H(y))=y\bigr)
  }{\rEE{3,4,17}}
\end{tabproofwide}

\subsubsection{Injektion aus Surjektion (via Sektion)}

% --- (1) Korollar: Surjektion + Sektion => Injektion ---
\FormulaThmDeltaR{%
  \exists H\colon B\inj A%
}{
 F\colon A\sur B\vdash \exists H\colon B\inj A
}{
  \DeltaRow{Mengen}{A\dsep B}
  \DeltaPrem{Surjektive Funktion}{F\colon A\sur B}
  \DeltaPrem{Funktion}{G\colon B\to A}
}
\begin{tabproof}
  \proofstep{1}{\exists S\colon B\to A\,\bigl(F\circ S=\Id_B\bigr)}{%
    \FormulaRefAuto{\exists F\colon A \to B\,G\circ F = \Id_A}}

  \proofstep{2}{S\colon B\to A \land F\circ S=\Id_B}{\rA}
  \proofstep{2}{F\circ S=\Id_B}{\rAEb{2}}

  \proofstep{2}{S\colon B\inj A}{%
    \FormulaRefAuto{F\circ G=\Id_B\vdash G\colon B\inj A}}

  \proofstep{2}{\exists H\colon B\inj A}{\rEI{4}}

  \proofstep{1}{\exists H\colon B\inj A}{\rEE{1,2,5}}
\end{tabproof}






\chapter{Faserweise Konstruktion von Bijektionen}

\ProofWideReasonsNearFormula

\section{Bijektionen zwischen Retraktionen}

Eine Retraktion \(r\colon S\to G\) auf eine Teilmenge \(G\subseteq S\)
fixiert jedes Element von \(G\). Ihre Faser über \(g\in G\) enthält daher
den ausgezeichneten Punkt \(g\). Bei zwei Retraktionen werden zunächst
Bijektionen der entsprechenden Fasern gewählt, die diese Punkte treffen.
Die Fasern zerlegen den Definitionsbereich und lassen sich anschließend
zu einer einzigen Bijektion zusammensetzen.

Wir verwenden im folgenden Satz die Abkürzungen
\[
  \begin{aligned}
    Q_g&=r^{-1}[\{g\}]=\{x\in S\mid r(x)=g\},\\[-2pt]
    R_h&=s^{-1}[\{h\}]=\{y\in T\mid s(y)=h\}.
  \end{aligned}
\]
Diese Gleichheiten folgen aus
\FormulaRefAuto{F\colon A\to B\dsep y\in B\vdash
  F^{-1}[\{y\}]=\{\,x\in A\mid F(x)=y\,\}}[thm]. Insbesondere gilt
\(x\in Q_g\leftrightarrow(x\in S\land r(x)=g)\), und entsprechend für
\(R_h\).

\FormulaThmDeltaK[Bijektion aus Bijektionen der Retraktionsfasern]{%
  \begin{aligned}[t]
    &\forall g\in G\,\exists u\,(u\colon Q_g\bij R_{\varphi(g)})
      \vdash{}\\[-2pt]
    &\exists f\,\Bigl(f\colon S\bij T
      \land\forall x\in S\,s(f(x))=\varphi(r(x))\\[-2pt]
    &\hspace{35mm}\land\forall g\in G\,f(g)=\varphi(g)\Bigr)
  \end{aligned}%
}{RetractionFiberBijection}{%
  \DeltaRow{Mengen}{S\dsep T\dsep G\dsep H\dsep
    g\dsep h\dsep x\dsep x'\dsep y\dsep z}
  \DeltaRow{}{Q_g\dsep R_h\dsep W\dsep D}
  \DeltaPrem{Teilmengen}{G\subseteq S\dsep H\subseteq T}
  \DeltaPrem{Funktionen}{r\colon S\to G\dsep s\colon T\to H\dsep
    \varphi\colon G\bij H}
  \DeltaPrem{Retraktionsgleichungen}{%
    \forall g\in G\,r(g)=g\dsep\forall h\in H\,s(h)=h}
  \DeltaRow{Hilfsfunktionen}{u\dsep U\dsep f\dsep\varphi^{-1}}
}
Im Beweis steht \(W=\powerset(S\times T)\) für einen gemeinsamen
Mengenbereich aller Faserbijektionen. Die Formel
\[
  \mathcal C(U)\;\coloneqq\;
  \forall g\in G\,
    \bigl(U(g)\colon Q_g\bij R_{\varphi(g)}
      \land U(g)(g)=\varphi(g)\bigr)
\]
beschreibt die gewünschten Werte einer Auswahl \(U\colon G\to W\).
Der folgende Verklebungsteil benötigt noch kein Auswahlprinzip.

Für die Surjektivität kürzt \(k_y\) den Term
\(\varphi^{-1}(s(y))\) ab.
\begin{tabproofsplitwide}
  \proofpartwide{Konstruktion und Funktionstyp}
    \proofstepwidestar[1]{U\colon G\to W}{%
      \rA}
    \proofstepwidestar[2]{\mathcal C(U)}{%
      \rA}
    \proofstepwidestar[3]{r\colon S\to G}{%
      \rA}
    \proofstepwidestar[4]{s\colon T\to H}{%
      \rA}
    \proofstepwidestar[5]{\varphi\colon G\bij H}{%
      \rA}
    \proofstepwidestar[6]{G\subseteq S}{%
      \rA}
    \proofstepwidestar[7]{\forall g\in G\,r(g)=g}{%
      \rA}
    \proofstepwidestar[5]{\varphi\colon G\to H}{%
      \FormulaRefAuto{Bijektive Funktion}[def]{5}}
    \proofstepwidestar[5]{\varphi\colon G\inj H}{%
      \FormulaRefAuto{Bijektive Funktion}[def]{5}}
    \proofstepwidestar[10]{x\in S}{%
      \rA}
    \proofstepwidestar[3,10]{r(x)\in G}{%
      \FormulaRefAuto{F\colon A\to B\dsep x\in A\vdash F(x)\in B}[thm]{3,10}}
    \proofstepwidestar[10]{x\in Q_{r(x)}}{%
      \FormulaRefAuto{x\in A\dsep P(x)\vdash x\in\{x\in A\mid P(x)\}}[thm]{10,\rII}}
    \proofstepwidestar[2,3,10]{U(r(x))\colon Q_{r(x)}\bij R_{\varphi(r(x))}}{%
      \rAEa{\rRE{\rUE{2},11}}}
    \proofstepwidestar[2,3,10]{U(r(x))\colon Q_{r(x)}\to R_{\varphi(r(x))}}{%
      \FormulaRefAuto{Bijektive Funktion}[def]{13}}
    \proofstepwidestar[2,3,10]{U(r(x))(x)\in R_{\varphi(r(x))}}{%
      \FormulaRefAuto{F\colon A\to B\dsep x\in A\vdash F(x)\in B}[thm]{14,12}}
    \proofstepwidestar[2,3,10]{U(r(x))(x)\in T\land s(U(r(x))(x))=\varphi(r(x))}{%
      \FormulaRefAuto{x\in\{u\in A\mid P(u)\}\eqvdash x\in A\land P(x)}[thm]{15}}
    \proofstepwidestar[2,3]{\forall x\in S\,U(r(x))(x)\in T}{%
      \rUI{\rRI{10,\rAEa{16}}}}
    \proofstepwidestar[2,3]{\exists f\,\bigl(f\colon S\to T\land\forall x\in S\,f(x)=U(r(x))(x)\bigr)}{%
      \FormulaRefAuto{TypedFunctionDefinitionPrinciple}[thm]{17}}
  \closeproofpartwide

  \proofpartwide{Verträglichkeit mit den Retraktionen}
    \setcounter{proofstepnr}{18}
    \proofstepwidestar[19]{f\colon S\to T\land\forall x\in S\,f(x)=U(r(x))(x)}{%
      \rA}
    \proofstepwidestar[19]{f\colon S\to T}{%
      \rAEa{19}}
    \proofstepwidestar[19]{\forall x\in S\,f(x)=U(r(x))(x)}{%
      \rAEb{19}}
    \proofstepwidestar[19,10]{f(x)=U(r(x))(x)}{%
      \rRE{\rUE{21},10}}
    \proofstepwidestar[2,3,19,10]{s(f(x))=\varphi(r(x))}{%
      \rIE{\rIE{22,\rII},\rAEb{16}}}
    \proofstepwidestar[2,3,19]{\forall x\in S\,s(f(x))=\varphi(r(x))}{%
      \rUI{\rRI{10,23}}}
  \closeproofpartwide

  \proofpartwide{Fortsetzung auf den Retraktionsbildern}
    \setcounter{proofstepnr}{24}
    \proofstepwidestar[25]{g\in G}{%
      \rA}
    \proofstepwidestar[6,25]{g\in S}{%
      \FormulaRefAuto{A\subseteq B\dsep x\in A\vdash x\in B}[thm]{6,25}}
    \proofstepwidestar[7,25]{r(g)=g}{%
      \rRE{\rUE{7},25}}
    \proofstepwidestar[2,25]{U(g)(g)=\varphi(g)}{%
      \rAEb{\rRE{\rUE{2},25}}}
    \proofstepwidestar[19,6,25]{f(g)=U(r(g))(g)}{%
      \rRE{\rUE{21},26}}
    \proofstepwidestar[2,19,6,7,25]{f(g)=\varphi(g)}{%
      \rIE{27,28,29}}
    \proofstepwidestar[2,19,6,7]{\forall g\in G\,f(g)=\varphi(g)}{%
      \rUI{\rRI{25,30}}}
  \closeproofpartwide

  \proofpartwide{Injektivität}
    \setcounter{proofstepnr}{31}
    \proofstepwidestar[32]{x'\in S}{%
      \rA}
    \proofstepwidestar[33]{f(x)=f(x')}{%
      \rA}
    \proofstepwidestar[3,32]{r(x')\in G}{%
      \FormulaRefAuto{F\colon A\to B\dsep x\in A\vdash F(x)\in B}[thm]{3,32}}
    \proofstepwidestar[2,3,19,32]{s(f(x'))=\varphi(r(x'))}{%
      \rRE{\rUE{24},32}}
    \proofstepwidestar[2,3,19,10,32,33]{\varphi(r(x))=\varphi(r(x'))}{%
      \rIE{23,35,\rIE{33,\rII}}}
    \proofstepwidestar[2,3,5,19,10,32,33]{r(x)=r(x')}{%
      \FormulaRefAuto{F\colon A\inj B\dsep x\in A\dsep y\in A\dsep F(x)=F(y)\vdash x=y}[ax]{9,11,34,36}}
    \proofstepwidestar[32]{x'\in Q_{r(x')}}{%
      \FormulaRefAuto{x\in A\dsep P(x)\vdash x\in\{x\in A\mid P(x)\}}[thm]{32,\rII}}
    \proofstepwidestar[2,3,5,19,10,32,33]{x'\in Q_{r(x)}}{%
      \rIE{\rIE{37,\rII},38}}
    \proofstepwidestar[19,32]{f(x')=U(r(x'))(x')}{%
      \rRE{\rUE{21},32}}
    \proofstepwidestar[2,3,5,19,10,32,33]{U(r(x))(x)=U(r(x))(x')}{%
      \rIE{22,40,\rIE{37,\rII},33}}
    \proofstepwidestar[2,3,10]{U(r(x))\colon Q_{r(x)}\inj R_{\varphi(r(x))}}{%
      \FormulaRefAuto{Bijektive Funktion}[def]{13}}
    \proofstepwidestar[2,3,5,19,10,32,33]{x=x'}{%
      \FormulaRefAuto{F\colon A\inj B\dsep x\in A\dsep y\in A\dsep F(x)=F(y)\vdash x=y}[ax]{42,12,39,41}}
    \proofstepwidestar[2,3,5,19]{\forall x\in S\,\forall x'\in S\,(f(x)=f(x')\rightarrow x=x')}{%
      \rUI{\rRI{10,\rUI{\rRI{32,\rRI{33,43}}}}}}
    \proofstepwidestar[2,3,5,19]{f\colon S\inj T}{%
      \FormulaRefAuto{Injektive Funktion}[def]{20,44}}
  \closeproofpartwide

  \proofpartwide{Surjektivität}
    \setcounter{proofstepnr}{45}
    \proofstepwidestar[46]{y\in T}{%
      \rA}
    \proofstepwidestar[4,46]{s(y)\in H}{%
      \FormulaRefAuto{F\colon A\to B\dsep x\in A\vdash F(x)\in B}[thm]{4,46}}
    \proofstepwidestar[5,4,46]{k_y\in G}{%
      \FormulaRefAuto{InverseFunctionValueType}[thm]{5,47}}
    \proofstepwidestar[5,4,46]{\varphi(k_y)=s(y)}{%
      \FormulaRefAuto{InverseFunctionRightInverse}[thm]{5,47}}
    \proofstepwidestar[5,4,46]{y\in R_{\varphi(k_y)}}{%
      \FormulaRefAuto{x\in A\dsep P(x)\vdash x\in\{x\in A\mid P(x)\}}[thm]{46,\rIE{49,\rII}}}
    \proofstepwidestar[2,5,4,46]{U(k_y)\colon Q_{k_y}\bij R_{\varphi(k_y)}}{%
      \rAEa{\rRE{\rUE{2},48}}}
    \proofstepwidestar[2,5,4,46]{U(k_y)\colon Q_{k_y}\sur R_{\varphi(k_y)}}{%
      \FormulaRefAuto{Bijektive Funktion}[def]{51}}
    \proofstepwidestar[2,5,4,46]{\exists z\in Q_{k_y}\,U(k_y)(z)=y}{%
      \FormulaRefAuto{Surjektivität}[ax]{52,50}}
    \proofstepwidestar[54]{z\in Q_{k_y}\land U(k_y)(z)=y}{%
      \rA}
    \proofstepwidestar[54]{z\in S\land r(z)=k_y}{%
      \FormulaRefAuto{x\in\{u\in A\mid P(u)\}\eqvdash x\in A\land P(x)}[thm]{\rAEa{54}}}
    \proofstepwidestar[19,54]{f(z)=U(r(z))(z)}{%
      \rRE{\rUE{21},\rAEa{55}}}
    \proofstepwidestar[19,54]{f(z)=y}{%
      \rIE{\rAEb{55},\rAEb{54},56}}
    \proofstepwidestar[19,54]{\exists z\in S\,f(z)=y}{%
      \rEI{\rAI{\rAEa{55},57}}}
    \proofstepwidestar[2,5,4,19,46]{\exists z\in S\,f(z)=y}{%
      \rEE{53,54,58}}
    \proofstepwidestar[2,5,4,19]{\forall y\in T\,\exists z\in S\,f(z)=y}{%
      \rUI{\rRI{46,59}}}
    \proofstepwidestar[2,5,4,19]{f\colon S\sur T}{%
      \FormulaRefAuto{Surjektive Funktion}[def]{20,60}}
  \closeproofpartwide

  \proofpartwideindR[Verkleben ausgewählter Faserbijektionen]{%
    U\colon G\to W\dsep\mathcal C(U)\vdash \begin{aligned}[t]&\exists f\,\bigl(f\colon S\bij T\land{}\\[-2pt]&\quad\forall x\in S\,s(f(x))=\varphi(r(x))\land{}\\[-2pt]&\quad\forall g\in G\,f(g)=\varphi(g)\bigr)\end{aligned}
  }{%
    U\colon G\to W\dsep\mathcal C(U)\vdash \begin{aligned}[t]&\exists f\,\bigl(f\colon S\bij T\land{}\\[-2pt]&\quad\forall x\in S\,s(f(x))=\varphi(r(x))\land{}\\[-2pt]&\quad\forall g\in G\,f(g)=\varphi(g)\bigr)\end{aligned}
  }[RetractionFiberBijectionGluing]
    \setcounter{proofstepnr}{61}
    \proofstepwidestar[2,3,4,5,19]{f\colon S\bij T}{%
      \FormulaRefAuto{F\colon A\inj B\dsep F\colon A\sur B\vdash F\colon A\bij B}[thm]{45,61}}
    \proofstepwidestar[2,3,4,5,6,7,19]{\begin{aligned}[t]&\exists f\,\bigl(f\colon S\bij T\land{}\\[-2pt]&\quad\forall x\in S\,s(f(x))=\varphi(r(x))\land{}\\[-2pt]&\quad\forall g\in G\,f(g)=\varphi(g)\bigr)\end{aligned}}{%
      \rEI{\rAI{62,\rAI{24,31}}}}
    \proofstepwidestar[2,3,4,5,6,7]{\begin{aligned}[t]&\exists f\,\bigl(f\colon S\bij T\land{}\\[-2pt]&\quad\forall x\in S\,s(f(x))=\varphi(r(x))\land{}\\[-2pt]&\quad\forall g\in G\,f(g)=\varphi(g)\bigr)\end{aligned}}{%
      \rEE{18,19,63}}
  \closeproofpartwideindR
\end{tabproofsplitwide}

Für den Auswahlteil setzen wir
\[
  D=\bigl\{(g,u)\in G\times W\bigm|
      u\colon Q_g\bij R_{\varphi(g)}\land u(g)=\varphi(g)\bigr\}.
\]
\begin{tabproofsplitwide}
  \setcounter{proofpartnr}{1}% Fortsetzung desselben Theorems nach dem Verklebungsteil.
  \proofpartwideindR[Totalität der Auswahlrelation]{%
    \forall g\in G\,\exists u\,(u\colon Q_g\bij R_{\varphi(g)})\vdash\TotRel{D,G,W}
  }{%
    \forall g\in G\,\exists u\,(u\colon Q_g\bij R_{\varphi(g)})\vdash\TotRel{D,G,W}
  }[RetractionFiberSelectionTotal]
    \proofstepwidestar[1]{\forall g\in G\,\exists u\,(u\colon Q_g\bij R_{\varphi(g)})}{%
      \rA}
    \proofstepwidestar[2]{r\colon S\to G}{%
      \rA}
    \proofstepwidestar[3]{s\colon T\to H}{%
      \rA}
    \proofstepwidestar[4]{\varphi\colon G\bij H}{%
      \rA}
    \proofstepwidestar[5]{G\subseteq S}{%
      \rA}
    \proofstepwidestar[6]{H\subseteq T}{%
      \rA}
    \proofstepwidestar[7]{\forall g\in G\,r(g)=g}{%
      \rA}
    \proofstepwidestar[8]{\forall h\in H\,s(h)=h}{%
      \rA}
    \proofstepwidestar[9]{g\in G}{%
      \rA}
    \proofstepwidestar[5,9]{g\in S}{%
      \FormulaRefAuto{A\subseteq B\dsep x\in A\vdash x\in B}[thm]{5,9}}
    \proofstepwidestar[7,9]{r(g)=g}{%
      \rRE{\rUE{7},9}}
    \proofstepwidestar[5,7,9]{g\in Q_g}{%
      \FormulaRefAuto{x\in A\dsep P(x)\vdash x\in\{x\in A\mid P(x)\}}[thm]{10,11}}
    \proofstepwidestar[4]{\varphi\colon G\to H}{%
      \FormulaRefAuto{Bijektive Funktion}[def]{4}}
    \proofstepwidestar[4,9]{\varphi(g)\in H}{%
      \FormulaRefAuto{F\colon A\to B\dsep x\in A\vdash F(x)\in B}[thm]{13,9}}
    \proofstepwidestar[4,6,9]{\varphi(g)\in T}{%
      \FormulaRefAuto{A\subseteq B\dsep x\in A\vdash x\in B}[thm]{6,14}}
    \proofstepwidestar[4,8,9]{s(\varphi(g))=\varphi(g)}{%
      \rRE{\rUE{8},14}}
    \proofstepwidestar[4,6,8,9]{\varphi(g)\in R_{\varphi(g)}}{%
      \FormulaRefAuto{x\in A\dsep P(x)\vdash x\in\{x\in A\mid P(x)\}}[thm]{15,16}}
    \proofstepwidestar[1,9]{\exists u\,(u\colon Q_g\bij R_{\varphi(g)})}{%
      \rRE{\rUE{1},9}}
    \proofstepwidestar[1,2,3,4,5,6,7,8,9]{\exists u\,\bigl(u\colon Q_g\bij R_{\varphi(g)}\land u(g)=\varphi(g)\bigr)}{%
      \FormulaRefAuto{PointedBijectionExists}[thm]{18,12,17}}
    \proofstepwidestar[20]{u\colon Q_g\bij R_{\varphi(g)}\land u(g)=\varphi(g)}{%
      \rA}
    \proofstepwidestar[20]{u\colon Q_g\to R_{\varphi(g)}}{%
      \FormulaRefAuto{Bijektive Funktion}[def]{\rAEa{20}}}
    \proofstepwidestar[20]{u\subseteq Q_g\times R_{\varphi(g)}}{%
      \FormulaRefAuto{FunctionGraphSubsetProduct}[thm]{21}}
    \proofstepwidestar[]{Q_g\subseteq S}{%
      \FormulaRefAuto{\{x\in A\mid P(x)\}\subseteq A}[thm]}
    \proofstepwidestar[]{R_{\varphi(g)}\subseteq T}{%
      \FormulaRefAuto{\{x\in A\mid P(x)\}\subseteq A}[thm]}
    \proofstepwidestar[]{Q_g\times R_{\varphi(g)}\subseteq S\times R_{\varphi(g)}}{%
      \FormulaRefAuto{CartesianProductMonotoneLeft}[thm]{23}}
    \proofstepwidestar[]{S\times R_{\varphi(g)}\subseteq S\times T}{%
      \FormulaRefAuto{CartesianProductMonotoneRight}[thm]{24}}
    \proofstepwidestar[20]{u\subseteq S\times R_{\varphi(g)}}{%
      \FormulaRefAuto{A\subseteq B\dsep B\subseteq C\vdash A\subseteq C}[thm]{22,25}}
    \proofstepwidestar[20]{u\subseteq S\times T}{%
      \FormulaRefAuto{A\subseteq B\dsep B\subseteq C\vdash A\subseteq C}[thm]{27,26}}
    \proofstepwidestar[20]{u\in W}{%
      \FormulaRefAuto{x\in\powerset(A)\;\eqvdash\;x\subseteq A}[thm]{28}}
    \proofstepwidestar[9,20]{(g,u)\in G\times W}{%
      \FormulaRefAuto{a\in A\dsep b\in B\vdash(a,b)\in A\times B}[thm]{9,29}}
    \proofstepwidestar[9,20]{(g,u)\in D}{%
      \FormulaRefAuto{x\in A\dsep P(x)\vdash x\in\{x\in A\mid P(x)\}}[thm]{30,20}}
    \proofstepwidestar[9,20]{\exists u\,((g,u)\in D)}{%
      \rEI{31}}
    \proofstepwidestar[1,2,3,4,5,6,7,8,9]{\exists u\,((g,u)\in D)}{%
      \rEE{19,20,32}}
    \proofstepwidestar[1,2,3,4,5,6,7,8]{\forall g\in G\,\exists u\,((g,u)\in D)}{%
      \rUI{\rRI{9,33}}}
    \proofstepwidestar[]{D\subseteq G\times W}{%
      \FormulaRefAuto{\{x\in A\mid P(x)\}\subseteq A}[thm]}
    \proofstepwidestar[1,2,3,4,5,6,7,8]{\TotRel{D,G,W}}{%
      \FormulaRefAuto{Totale Relation}[def]{35,34}}
  \closeproofpartwideindR

  \proofpartwide{Auswahl und Schluss}
    \proofstepwidestar[1]{\forall g\in G\,\exists u\,(u\colon Q_g\bij R_{\varphi(g)})}{%
      \rA}
    \proofstepwidestar[1]{\TotRel{D,G,W}}{%
      \FormulaRefAuto{RetractionFiberSelectionTotal}[thm]{1}}
    \proofstepwidestar[1]{\exists U\,\bigl(U\colon G\to W\land\forall g\in G\,(g,U(g))\in D\bigr)}{%
      \FormulaRefAuto{\exists F\colon A\to B\,\forall x\in A\,\bigl((x,F(x))\in R\bigr)}[thm]{2}}
    \proofstepwidestar[4]{U\colon G\to W\land\forall g\in G\,(g,U(g))\in D}{%
      \rA}
    \proofstepwidestar[4]{U\colon G\to W}{%
      \rAEa{4}}
    \proofstepwidestar[6]{g\in G}{%
      \rA}
    \proofstepwidestar[4,6]{(g,U(g))\in D}{%
      \rRE{\rUE{\rAEb{4}},6}}
    \proofstepwidestar[4,6]{(g,U(g))\in G\times W\land\bigl(U(g)\colon Q_g\bij R_{\varphi(g)}\land U(g)(g)=\varphi(g)\bigr)}{%
      \FormulaRefAuto{x\in\{u\in A\mid P(u)\}\eqvdash x\in A\land P(x)}[thm]{7}}
    \proofstepwidestar[4]{\mathcal C(U)}{%
      \rUI{\rRI{6,\rAEb{8}}}}
    \proofstepwidestar[4]{\begin{aligned}[t]&\exists f\,\bigl(f\colon S\bij T\land{}\\[-2pt]&\quad\forall x\in S\,s(f(x))=\varphi(r(x))\land{}\\[-2pt]&\quad\forall g\in G\,f(g)=\varphi(g)\bigr)\end{aligned}}{%
      \FormulaRefAuto{RetractionFiberBijectionGluing}[thm]{5,9}}
    \proofstepwidestar[1]{\begin{aligned}[t]&\exists f\,\bigl(f\colon S\bij T\land{}\\[-2pt]&\quad\forall x\in S\,s(f(x))=\varphi(r(x))\land{}\\[-2pt]&\quad\forall g\in G\,f(g)=\varphi(g)\bigr)\end{aligned}}{%
      \rEE{3,4,10}}
  \closeproofpartwide
\end{tabproofsplitwide}

Die beiden punktweisen Schlussgleichungen bedeuten nach
\FormulaRefAuto{FunctionExtensionality}[thm] auch
\(s\circ f=\varphi\circ r\) und \(f\restriction_G=\varphi\).
Nur die gemeinsame Auswahl der Faserbijektionen benutzt hier AC.

\ProofWideReasonsAtRight

\end{document}
